\documentclass[12pt]{article}

% ── Packages ──────────────────────────────────────────────────────────────────
\usepackage[margin=1in]{geometry}
\usepackage{setspace}
\usepackage{times}
\usepackage[T1]{fontenc}
\usepackage[utf8]{inputenc}
\usepackage[style=apa, backend=biber]{biblatex}
\usepackage{hyperref}
\usepackage{amsmath}
\usepackage{booktabs}
\usepackage{graphicx}
\usepackage{microtype}
\usepackage{indentfirst}
\usepackage{parskip}

\addbibresource{references.bib}

% ── Layout ────────────────────────────────────────────────────────────────────
\doublespacing
\setlength{\parindent}{0.5in}
\setlength{\parskip}{0pt}

\hypersetup{
  colorlinks = true,
  linkcolor  = black,
  citecolor  = black,
  urlcolor   = blue
}

% ══════════════════════════════════════════════════════════════════════════════
\begin{document}

% ── Title Page ────────────────────────────────────────────────────────────────
% FIX 1 & 2: Title updated to match paper content (removed "Home Pool Advantage"
% framing since home pool advantage is not directly analyzed in the results).
\begin{titlepage}
  \centering
  \vspace*{2in}

  {\large\bfseries Pool Advantage: Facility Characteristics and Performance
  in NCAA Swimming\par}
  \vspace{1cm}

  {\large Lily Cramer\par}
  \vspace{0.3cm}
  {\large Department of Sport Management, Rice University\par}
  \vspace{0.3cm}
  {\large SMGT 490: Dr.\ Hua Gong\par}
  \vspace{0.3cm}
  {\large May 3, 2026\par}

  \vfill
\end{titlepage}

\setcounter{page}{1}

% ── Table of Contents ─────────────────────────────────────────────────────────
\tableofcontents
\newpage

% ══════════════════════════════════════════════════════════════════════════════
\section{Introduction}

Athletes routinely characterize venues as ``fast'' or ``slow.'' This judgment
is passed almost immediately upon entry, based on the feel of the water, the
energy of the environment, and most concretely the time posted on the
scoreboard. A swimmer who drops several seconds at one facility but gains time
at another may attribute the difference to the pool itself rather than to their
own preparation. This shared vocabulary, so embedded in the culture of
swimming, raises a straightforward empirical question: is there anything to it?

Unlike professional sports leagues such as the NBA or NFL, which operate under
strictly standardized facility requirements, the NCAA imposes no uniform
specifications beyond pool length on the venues at which its athletes compete.
Lane width, water depth, gutter design, lighting, water temperature, and
altitude can all vary substantially from one facility to the next. Some pools
are housed in dated natatoriums with limited spectator capacity; others are
purpose-built competition facilities engineered for elite performance. These
differences are not solely cosmetic.

\subsection{Research Motivation}

Despite the pervasiveness of ``fast pool'' discourse among athletes and
coaches, surprisingly little research has attempted to quantify the effect of
venue characteristics on swimming performance. The existing sports analytics
literature has devoted considerable attention to home-field advantage and
stadium effects across a range of professional sports, but competitive swimming
has received comparatively little scholarly attention. The anecdotal evidence
is robust: coaches plan travel schedules around prestigious fast pools, and
athletes describe palpable differences in how they feel and perform across
facilities. However, no large-scale quantitative study has validated these
observations in the context of NCAA competition.

This gap in the literature is the primary motivation for the present study. As
a competitive swimmer, the author has experienced firsthand the apparent
influence of pool environment on performance and has observed teammates and
competitors respond similarly across venues. The goal of this research is to
move beyond anecdote and subject the ``fast pool'' hypothesis to rigorous
empirical scrutiny.

\subsection{Research Questions}

This study is organized around one primary research question and two supporting
secondary questions. The primary question asks: what impact does a pool have on
a swimmer's performance in NCAA competition? Specifically, to what extent do
times achieved at a given venue systematically deviate from what would be
predicted based on athlete ability alone?

The secondary questions address two related dimensions. First, is there
evidence of a general ``pool advantage,'' meaning a venue effect that persists
after controlling for the quality of athletes competing at each facility?
Second, which structural or environmental characteristics of a pool are most
strongly associated with faster times? Candidate factors include water depth,
lane configuration, altitude, indoor versus outdoor environment, spectator
capacity, and facility age.

\subsection{Contribution Overview}

This paper makes several contributions to the sports analytics literature. To
the author's knowledge, it presents the first large-scale empirical study of
pool effects in NCAA swimming, drawing on performance data across a broad
sample of athletes, events, and venues. In doing so, it extends the growing
body of research on environmental and contextual factors in athletic performance
into a sport that has been largely overlooked by quantitative researchers.

Beyond its scholarly contribution, this study has practical implications for
athletes, coaches, and meet organizers. If certain pool characteristics are
reliably associated with faster times, this information can inform decisions
about where to schedule high-stakes competitions, how to interpret times posted
at different facilities, and what investments in facility design might yield the
greatest performance dividends. At a moment when data-driven decision-making is
reshaping nearly every corner of competitive sport, swimming stands to benefit
from a more rigorous understanding of how its physical environments shape what
happens in the water.

% ══════════════════════════════════════════════════════════════════════════════
\section{Literature Review}

\subsection{Structural and Regulatory Influences on Performance}

Competitive swimming facilities are subject to regulatory oversight designed to
ensure consistency across venues. World Aquatics mandates standards for pool
construction, including a length tolerance of no more than 0.01 meters beyond
the standard 50-meter distance, minimum water temperatures of 25 degrees
Celsius, and current limits of 1.25 meters per minute \parencite{fina2021}.
% FIX 6: Menting et al. (2019) is a pacing paper, not a lane line paper.
% Citation replaced with the FINA source which covers equipment standards,
% and the claim is reframed accordingly.
Lane line specifications intended to reduce turbulence from adjacent swimmers
are likewise governed by World Aquatics equipment rules, though compliance
varies across NCAA venues \parencite{fina2021}.

Despite these standardization efforts, performance outcomes remain sensitive to
the physical environment of the pool. Competitive swimming races consist of
multiple phases, including the start, underwater breakout, open-water swimming,
turns, and finish, each of which is affected differently by facility design.
Research has shown that even small variations within regulatory guidelines can
have a measurable impact on swimmer performance across these phases
\parencite{marinho2020, morais2019, nicol2021, veiga2016}. \textcite{born2022}
further establish that turn performance carries the largest effect on race time
in short-course events, with a standardized coefficient of 0.53 or greater,
making it a particularly important phase for any model that aims to capture
pool-level influences on performance.

\subsection{Environmental Conditions and Performance}

Environmental conditions within aquatic facilities extend beyond structural
dimensions and play a significant role in shaping both physical capacity and
psychological readiness. Competing or training at altitude induces physiological
adaptations including increased red blood cell production and enhanced aerobic
endurance, though hypoxic conditions may also alter markers such as blood
lactate concentration \parencite{perrimoore2015}. Indoor air quality and
temperature represent additional environmental variables of relevance: warmer
ambient temperatures and exposure to air pollutants in enclosed natatoriums have
been associated with reduced endurance and impaired cognitive function
\parencite{sargeant2016}.

Psychological responses to the facility environment have also been shown to
influence athletic performance. Athletes competing in facilities that are
well-maintained and aesthetically designed report higher levels of confidence
and motivation \parencite{alasinrin2024, lynch2018}. Facilities incorporating
natural lighting have been linked to improved mood and cognitive functioning,
while reliance on artificial lighting alone has been associated with less
favorable psychological outcomes \parencite{scanlan2016}. These findings
collectively suggest that the facility environment shapes performance through
both physiological and psychological channels.

\subsection{Biomechanical Efficiency and Pool Design}

Swimming performance is uniquely sensitive to wave resistance and drag, and
pool design interacts directly with these forces. Underwater swimming following
starts and turns produces substantially less resistance than surface swimming,
and elite swimmers strategically maximize underwater distance to exploit this
advantage \parencite{vernell2006}. Research demonstrates that international-level
swimmers cover greater underwater distances per race than national-level
swimmers, reflecting more efficient exploitation of pool-specific hydrodynamic
conditions \parencite{pla2021}. Underwater propulsion is driven primarily by the
dolphin kick, which depends on coordinated undulation as well as ankle strength
and flexibility \parencite{shimojo2019, willems2014}. Pool characteristics that
affect turbulence and wave reflection therefore have direct implications for the
efficiency of this phase of the race.

\subsection{Methodological Considerations: Mixed-Effects Models in Swimming
Research}

The statistical framework best suited to estimating pool effects in this setting
is the mixed-effects regression model. \textcite{avalos2003} argue that mixed
models are particularly well suited to datasets involving repeated measurements
across multiple subjects, as they accommodate both intra-individual variation
across observations and inter-individual variation across athletes. This is
directly applicable to the present study, in which each swimmer contributes
multiple performances across different venues and events.

A key advantage of the mixed-effects framework is its ability to separate fixed
environmental influences from the random performance trajectories of individual
athletes. \textcite{baker2023} demonstrate this property in the context of
age-dependent career performance modeling, showing that mixed models effectively
capture autocorrelation in athletic performance and allow for individual-specific
% FIX 7: Added one sentence acknowledging the Baker & McHale context differs
% from the present study (golf vs. swimming) to pre-empt reviewer skepticism.
deviations from population-level trends. Although their application is to golf
rather than swimming, the statistical principles transfer directly: the same
partial-pooling logic that separates aging effects from cohort differences in
their framework is used here to separate pool effects from team and athlete
quality. Their work also highlights the risk of survivor bias in performance
datasets, wherein athletes still competing at advanced stages of their careers
appear to outperform historical norms simply because lower-ability athletes have
exited the sample. This concern is relevant to the present analysis, where pools
hosting competition late in the season may attract a positively selected sample
of athletes, and it motivates the inclusion of meet-level controls in the
primary specification.

\subsection{Research Gap}

Collectively, the literature demonstrates that swimming performance is shaped by
a complex combination of structural, environmental, biomechanical, and
psychological factors. These elements have been examined in relative isolation,
but no existing study has evaluated whether their combined influence produces a
measurable and generalizable pool advantage in the context of NCAA competition.
Given that swimmers typically train and compete within a single primary
facility, repeated exposure to specific environmental and structural conditions
may produce performance adaptations that persist across venues. Addressing this
gap has practical implications for facility selection, performance evaluation,
and competitive equity at the collegiate level.

% ══════════════════════════════════════════════════════════════════════════════
\section{Data}
\label{sec:data}

\subsection{Data Sources}

\subsubsection{USA Swimming}

The primary source of performance data for this study is the USA Swimming
database, which publishes the top 500 recorded times for each event in women's
NCAA swimming on a rolling basis throughout the competitive season
\parencite{usaswimming2026}. Data was collected for the 2025--2026 season,
capturing the fastest performances across all stroke-and-distance combinations
contested in collegiate competition. For each observation, the dataset provides
the swimmer's name, the event contested, the time recorded, the team
affiliation, and the meet at which the performance was achieved.

Because this source is limited to the top 500 times per event, the analytical
sample reflects the upper tier of NCAA Division I women's swimming rather than
the full population of collegiate performers, a feature of the data that is
relevant to the interpretation of results and is discussed further in
Section~\ref{sec:results}. The validity and utility of top-times databases as
a tool for performance evaluation in elite swimming has been established in
prior methodological work \parencite{smith2002}.

\subsubsection{Facility Data Collection}

Pool-level structural and environmental data were collected directly from the
official websites of each aquatic center represented in the sample. For each
facility, publicly available information was obtained on pool depth, lane count,
spectator capacity, facility age, and whether the pool is housed indoors or
outdoors. Where a facility's website provided incomplete information,
supplementary sources including university athletic department pages and
conference meet documentation were consulted. The depth and consistency of
facility documentation varies considerably across institutions, and this
variability introduces measurement uncertainty into several structural
covariates, as discussed in the limitations section.

\subsection{Sample Construction}

The analytical sample was constructed in two stages. In the first stage, all
individual-event observations from the USA Swimming top 500 lists for the
2025--2026 women's NCAA season were compiled across the 13 contested events.
Championship meets --- including conference championships and NCAA postseason
competition --- were excluded from the analytical sample at this stage. This
restriction is motivated by selection bias concerns: championship meets
disproportionately attract the most elite swimmers to a small number of
high-profile venues, meaning that pools hosting postseason competition would
appear artificially fast simply because of the quality of athletes competing
there rather than any genuine venue effect. Restricting the sample to standard
invitational meets ensures that the variation in performance across pools more
credibly reflects venue characteristics rather than the self-selection of elite
athletes into particular facilities.

From this collection of standard-meet performances, the 28 pools with the
greatest number of entries in the dataset were identified and retained for
analysis. This threshold was selected to ensure that pool-level effect estimates
are based on a sufficient number of observations to support reliable inference,
and to focus the facility data collection effort on venues with meaningful
representation in the performance data.

In the second stage, structural and environmental characteristics were collected
for each of the 28 pools from their respective aquatic center websites. Pools
for which sufficient facility documentation could not be obtained were flagged,
and analyses involving specific structural covariates are restricted
accordingly. The resulting dataset joins individual performance observations
from the top 500 lists with pool-level characteristics for the 28 most
represented venues in women's NCAA swimming during the 2025--2026 season.

Several features of this sample construction warrant acknowledgment. Because
the source data are limited to the top 500 times per event, slower programs and
lower-profile meets are underrepresented even within the standard-meet
subsample. The 28-pool restriction also concentrates the sample on high-traffic
venues, which tend to be larger and better-resourced facilities that host
multiple meets per season. These features affect the generalizability of the
findings and are discussed in the limitations section.

\subsection{Key Variables}

\subsubsection{Outcome Variable}

The primary outcome of interest is a swimmer's standardized race time within a
given stroke and distance combination. Raw race times are not directly
comparable across events, as a 200 butterfly time and a 50 freestyle time
occupy entirely different scales, so standardization is necessary to pool
observations across events in any analysis. For each stroke and distance
category, times are standardized by subtracting the sample mean and dividing by
the sample standard deviation, yielding a $z$-score that reflects how a given
swim compares to the broader distribution of performances in that event. Lower
standardized times indicate a faster performance. In event-specific models, raw
times are used directly.

\subsubsection{Primary Explanatory Variables}

% FIX 3 & 5: Added indoor/outdoor as a primary explanatory variable, which was
% missing from this section despite being the most significant result.
% Spectator capacity retained here but explicitly noted as dropped from the
% scaled model in the Results section.
The central explanatory variables capture pool-level structural and
environmental characteristics that hypothetically influence performance. Pool
depth is included given arguments that deeper pools reduce wave reflection from
the bottom, in turn decreasing turbulence. Altitude is included as a continuous
variable reflecting the elevation of the city in which the facility is housed.
Whether a pool is indoors or outdoors is included as a binary indicator,
reflecting the hypothesis that enclosed environments stabilize water temperature,
reduce wind interference, and eliminate outdoor environmental variability.
Facility age serves as a proxy for gutter and lane design, on the premise that
pools constructed more recently may differ from older facilities in their
wave-reduction profiles. Spectator capacity is included as a continuous variable
reflecting facility scale and the potential influence of crowd size on
performance, though it is omitted from the primary scaled model due to
multicollinearity with other facility-level covariates.

\subsubsection{Control Variables}

Several variables are included as controls to isolate the pool effect from
confounding sources. Stroke and distance are included to account for the
event-specific nature of the sport. Because championship meets have been
excluded from the sample, meet-level variation is addressed at the sample
construction stage rather than through a model control; this approach is
preferred over including meet level as a covariate because it removes the
source of selection bias entirely rather than attempting to adjust for it
post-hoc. Season timing captures where in the competitive calendar a
performance occurs, reflecting the expectation that athletes peak toward the
end of the season.

\subsubsection{Random Effects Structure}

The hierarchical structure of the data, with swimmers nested within teams and
performances nested within pools and meets, inspires a mixed-effects modeling
framework. Team-level random intercepts are included to account for differences
in the caliber of athletes across programs, ensuring that pools that
disproportionately host strong teams are not automatically identified as fast
venues. Pool-level clustering is addressed through pool random intercepts in the
full mixed model and through standard error clustering in robustness checks.
This structure allows the analysis to separate the contribution of the pool
environment from the influence of athlete quality, team strength, and
competitive context.

% ══════════════════════════════════════════════════════════════════════════════
\section{Exploratory Data Analysis}

Before estimating pool effects, this study conducts exploratory data analysis
to characterize the distribution of the outcome variable, assess cross-pool
variation in facility characteristics, and evaluate the feasibility of the
identification strategy employed in later models.

\subsection{Standardization of Race Times}

Because raw swim times vary substantially across events---spanning roughly 20
seconds in the 50 freestyle to over 15 minutes in the 1650 freestyle---direct
comparison of performances across event categories is not meaningful without
some transformation. Times are therefore standardized within each stroke and
distance combination using $z$-scores, centering each observation on the event
mean and scaling by the event standard deviation, consistent with
Section~\ref{sec:data}. Inspection of the resulting distributions confirms that
all 13 events are approximately centered at zero, validating the standardization
procedure. Each distribution exhibits a left skew, a pattern consistent with
the competitive structure of NCAA swimming: elite performances cluster near a
lower bound, while the upper tail of slower times exhibits greater spread. This
asymmetry is present across events and is accounted for in model specification.

\subsection{Unadjusted Variation in Performance Across Facilities}

% FIX 8: Removed the implication that a figure exists here since no figure
% is shown for this section. Reframed as "examination of" rather than "plotting."
Examination of median standardized times by facility reveals meaningful
unadjusted variation across pools. Several venues, including McAuley Aquatic
Center and the UVA Aquatic and Fitness Center, are associated with below-average
standardized times in the raw data, while others, such as the Florida Aquatics
Swimming and Training facility, are associated with above-average times. These
unadjusted differences may suggest a pool effect but must be interpreted with
caution. The interquartile ranges associated with most facilities are wide,
reflecting substantial heterogeneity in the caliber of athletes competing at
each venue. Pools that host elite invitational meets or conference championships
will inherently attract faster swimmers, making it difficult to attribute raw
performance differences to the pool itself. This pattern motivates the need for
a model that explicitly controls for athlete and team-level ability before
drawing conclusions about venue effects.

\subsection{Variation in Facility Characteristics}

Examination of the facility-level covariates confirms sufficient cross-pool
variation to support multivariate regression analysis. Altitude ranges from near
sea level to approximately 1,000 feet above sea level. Spectator capacity spans
roughly 400 to 2,500 seats. Pool depth ranges from approximately 90 to 150
inches across sampled facilities. The sampled pools are predominantly indoor
facilities, though a small number of outdoor venues are represented, providing
sufficient variation to identify an indoor effect. Lane count is nearly uniform
across Division I facilities and is unlikely to function as a meaningful
predictor. Facility age varies substantially, with some pools constructed within
the past decade and others exceeding 40 years old, a range sufficient to capture
meaningful generational differences in gutter design, lane rope quality, and
overall construction.

\subsection{Performance Variation by Meet Level}

Because championship meets were excluded from the analytical sample to address
selection bias, all performances in the dataset reflect standard invitational
competition. Examination of performance distributions across the remaining meet
types confirms that times are broadly comparable in competitive intensity,
supporting the validity of pooling standard-meet observations for the purpose
of estimating pool effects. This restriction is conservative by design: any
residual variation in athlete quality across standard meets is addressed through
the team-level random effects structure described in Section~\ref{sec:data}.

% ══════════════════════════════════════════════════════════════════════════════
\section{Empirical Strategy}

\subsection{Overview}

The central challenge in estimating pool effects is separating the influence of
venue characteristics from the influence of the athletes competing in them.
Pools that host high-profile competition attract stronger programs and more
elite swimmers, meaning that raw performance differences across facilities
conflate venue quality with athlete quality. Any credible estimate of a pool
effect must account for the systematic variation in competitor ability across
venues. The empirical strategy adopted here addresses this challenge through a
mixed-effects regression framework that controls for team strength, handles the
nested and repeated-measures structure of the data, and isolates venue-level
variation in performance after conditioning on observable confounders.

\subsection{Primary Model: Mixed-Effects Regression}

The primary model regresses standardized race time on a vector of pool-level
characteristics and a set of race-level controls, with a random intercept
specified at the team level:

\begin{equation}
  \text{RaceTime}_{ijkt} = \beta_0 + \beta_1 \text{PoolChar}_k +
  \beta_2 \text{Stroke}_j + \beta_3 \text{Distance}_j +
  u_i + \varepsilon_{ijkt}
  \label{eq:primary}
\end{equation}

\noindent where $i$ indexes teams, $j$ indexes events, $k$ indexes pools, and
$t$ indexes meets. The term $u_i$ represents a team-level random intercept
assumed to follow a normal distribution with mean zero and estimated variance,
and $\varepsilon_{ijkt}$ is the observation-level residual.

The pool characteristics vector includes the five facility-level covariates
described in Section~\ref{sec:data}: pool depth, altitude, indoor/outdoor
classification, facility age, and lane count. Race-level controls for stroke
and distance absorb the systematic variation in performance across event
categories that standardization does not fully eliminate, particularly
differences in the number of turns, which varies across distances and affects
the proportion of race time spent in the underwater phase. Meet level is not
included as a model covariate because championship meets have been excluded
from the sample entirely at the data construction stage, removing the primary
source of meet-level selection bias rather than attempting to adjust for it
within the model.

The mixed-effects specification is justified on statistical and substantive
grounds. Statistically, the data exhibit a hierarchical structure in which
swimmers are nested within teams and performances are nested within pools and
meets, a structure that violates the independence assumption of ordinary least
squares regression. Ignoring this structure would produce downward-biased
standard errors and artificially precise inference. Substantively, team-level
random intercepts absorb stable differences in program quality, ensuring that
pools disproportionately visited by strong programs are not spuriously
identified as fast venues.

\subsection{Identification Strategy}

The identification of pool effects rests on within-sample variation in the
pools at which a given team competes. Because many programs travel to multiple
venues across a season---including home meets and neutral-site
invitationals---the data contains observations of the same team performing at
different pools under varying facility conditions. Controlling for team-level
ability through the random intercept, the remaining variation in performance
across venues can be attributed to pool-level factors rather than to differences
in who is swimming.

The exclusion of championship meets strengthens this identification strategy
considerably. Championship pools are not randomly selected --- they tend to be
larger, newer, and better-resourced facilities that would appear fast regardless
of their structural characteristics simply because the best swimmers in the
country compete there. By restricting the sample to standard invitationals, the
remaining cross-pool variation in performance is more plausibly attributable to
venue characteristics than to the selection of athletes into meets.

\subsection{Robustness Checks}

Several alternative specifications are estimated to assess the sensitivity of
the primary results. First, pool fixed effects are substituted for the
continuous facility characteristic variables, replacing the parametric
assumptions of the primary model with a fully flexible set of pool-level
intercepts. A lower AIC in the fixed effects model would indicate that the
measured characteristics fail to fully capture the dimensions along which pools
differ.

Second, the model is re-estimated separately by stroke category to assess
whether pool characteristic effects are consistent across event types or
concentrated in particular disciplines. Third, alternative standardization
approaches for the outcome variable are evaluated, including log transformation
of raw times and percentile-based normalization, to confirm that results are not
sensitive to the specific transformation applied.

% ══════════════════════════════════════════════════════════════════════════════
\section{Results}
\label{sec:results}

\subsection{Model Assumptions and Diagnostics}

Prior to interpreting the primary model estimates, the mixed-effects
specification was assessed against standard regression assumptions. Examination
of the residuals versus fitted values plot reveals a modest fan shape, wherein
variance is somewhat larger among observations with low fitted values, that is,
among the fastest swims. This heteroskedasticity is consistent with the
competitive structure of NCAA swimming, where championship meets attract a wider
range of athlete abilities, producing greater dispersion in performance at the
elite end of the distribution. While this pattern represents a technical
violation of the homoskedasticity assumption, it does not materially compromise
inference given the large sample size and is acknowledged as a limitation of
the current specification.

The normal Q-Q plot indicates that residuals follow the theoretical normal
distribution closely across the central range of the data. Meaningful deviation
occurs only in the left tail, corresponding to exceptionally fast outlier
performances. This pattern mirrors the left skew observed in the exploratory
data analysis and reflects the physiological floor on swim times rather than a
modeling failure.

\subsection{Team Random Effects}

The estimated team random effects reveal substantial heterogeneity in baseline
performance across programs, with intercepts ranging from approximately $-1.0$
standard deviations below average for programs such as Virginia, Tennessee,
Texas, and Stanford, to approximately $+0.5$ standard deviations above average
for programs such as Ball State, Illinois, and Florida State. This wide
dispersion confirms that team strength varies enormously across the sample and
validates the decision to model team as a random effect. Had team ability been
omitted from the specification, its influence would have been absorbed into the
pool characteristic coefficients, producing substantially biased estimates of
venue effects.

\subsection{Primary Model Results}

% FIX 1: Figure cross-references now use \ref{} correctly so they resolve
% in the compiled PDF rather than showing "??".
Figure~\ref{fig:coef} presents the fixed effect estimates from the primary
mixed-effects model for each pool characteristic, with 95\% confidence intervals
displayed as horizontal bars. Points to the left of zero indicate
characteristics associated with faster swim times; points to the right indicate
characteristics associated with slower times. Gray points represent effects that
are not statistically distinguishable from zero, meaning their confidence
intervals cross the zero line. Blue points represent statistically significant
effects at the $p < 0.05$ threshold.

Of the five facility characteristics included, three emerge as statistically
significant predictors of standardized swim time. Being an indoor pool is the
largest and most precisely estimated effect in the scaled model, associated with
a reduction of approximately 0.24 standard deviations in swim time. This
finding is consistent with the hypothesis that enclosed environments stabilize
water temperature, reduce wind interference, and eliminate outdoor environmental
variability that might otherwise impair performance. Altitude also exhibits a
significant negative association with swim time, with higher-elevation
facilities associated with modestly faster performances. While the physiological
mechanism---reduced air resistance and aerobic adaptation to hypoxic
training---is well-documented in the literature, the relatively limited altitude
range in this sample constrains the magnitude of the observed effect. Facility
age is negatively and significantly associated with swim time as well, indicating
that older pools are associated with marginally faster performances. One
plausible explanation is survivor bias: pools that continue to host NCAA
competition decades after their construction tend to be well-established,
high-profile facilities that have been maintained and upgraded over time, rather
than a representative sample of aging infrastructure.

Pool depth does not reach statistical significance in the primary model, with a
confidence interval that crosses zero. While deeper pools are theoretically
associated with reduced wave reflection from the pool floor, the limited
variation in depth across Division I facilities may constrain the statistical
power available to detect this effect. Lane count is similarly non-significant;
as noted in the exploratory analysis, lane count is nearly uniform across the
sampled venues, making it difficult for the model to isolate its independent
contribution to performance. The wide confidence interval around the lane count
estimate reflects this limited identifying variation and should not be
interpreted as evidence that lane configuration is unimportant in general.

\begin{figure}[h]
  \centering
  % Export your R coefficient plot and upload to Overleaf, then replace the
  % placeholder below with:
  % \includegraphics[width=\textwidth]{Figures/coef_plot.png}
  \fbox{\parbox{0.85\textwidth}{\centering\vspace{2.5cm}
    [Insert Figure 1: Coefficient plot --- Effect of Facility Characteristics
    on Swim Performance]
  \vspace{2.5cm}}}
  \caption{Effect of Facility Characteristics on Swim Performance. Points to
  the left of zero are associated with faster swim times. Gray = confidence
  interval crosses zero (not significant at $p < 0.05$). Bars show 95\%
  confidence intervals.}
  \label{fig:coef}
\end{figure}

\subsection{Case Study: Christiansburg Aquatic Center}

To illustrate how facility characteristics combine to produce a pool-level
speed advantage, Figure~\ref{fig:contribution} decomposes the estimated
performance contribution of each structural characteristic at Christiansburg
Aquatic Center, the fastest pool in the sample after controlling for athlete and
team quality. Each bar represents the product of the pool's characteristic value
and the corresponding model coefficient, reflecting how much that feature
contributes to faster or slower swim times relative to an average facility.

The decomposition reveals that Christiansburg's speed advantage is driven
primarily by two characteristics: its indoor environment (contributing
$-0.245$ SD, the largest single factor) and its altitude (contributing
$-0.203$ SD). Together, these two features account for nearly half a standard
deviation of performance improvement relative to an average pool. Partially
offsetting these gains, pool depth works against performance at this facility
(contributing $+0.105$ SD), suggesting the pool is shallower than average.
Facility age adds a modest countervailing contribution ($+0.036$ SD), and lane
count is effectively neutral ($+0.006$ SD), consistent with its non-significant
coefficient in the primary model.

This decomposition illustrates a broader finding of the analysis: no single
characteristic defines a fast pool. Rather, venue speed emerges from a
combination of features that may partially offset one another, making simple
rankings of facilities by any one structural dimension an incomplete basis for
performance evaluation.

\begin{figure}[h]
  \centering
  % \includegraphics[width=\textwidth]{Figures/Characteristics.png}
  \fbox{\parbox{0.85\textwidth}{\centering\vspace{2.5cm}
    [Insert Figure 2: Christiansburg contribution plot]
  \vspace{2.5cm}}}
  \caption{Why is Christiansburg Aquatic Center the Fastest Pool? Each bar
  represents that characteristic's contribution to standardized swim time
  (model coefficient $\times$ pool value). Blue bars indicate characteristics
  associated with faster times; red bars indicate characteristics associated
  with slower times.}
  \label{fig:contribution}
\end{figure}

\subsection{Model Comparison and Robustness}

To assess whether the five measured pool characteristics adequately capture
venue-level variation, the primary characteristics model was compared against an
alternative specification replacing facility variables with pool fixed effects.
The fixed effects model produces a lower AIC (10,538.97 vs.\ 10,563.77) despite
the additional parameters required, indicating superior fit to the data. This
result suggests that NCAA pools differ from one another in ways the measured
structural and environmental variables do not fully capture. Unmeasured factors
such as gutter design, water temperature, lighting conditions, or the
psychological familiarity of a home venue may account for residual venue-level
variation that the characteristics model leaves unexplained. This finding
motivates future work to collect more granular facility data, as discussed in
Section~\ref{sec:limitations}.

% ══════════════════════════════════════════════════════════════════════════════
\section{Data Limitations}
\label{sec:limitations}

No dataset of this scope is without limitations. The limitations fall into
three broad categories: structural reporting inconsistencies, measurement
constraints, and sample size imbalances.

\subsection{Structural Reporting Inconsistencies}

The most significant data challenge concerns the inconsistent reporting of
facility characteristics across NCAA institutions. Pool depth, gutter
configuration, and lane line specifications are not systematically documented
in any centralized source. University athletic websites vary considerably in the
technical detail they provide about their natatoriums, and many offer no
structural specifications whatsoever. As a result, these variables are missing
for a non-trivial share of facilities in the sample, and the missingness is
unlikely to be random. Older, lower-profile, or under-resourced programs may be
less likely to publish detailed facility documentation, introducing the
possibility that missingness is correlated with both facility quality and
competitive outcomes.

\subsection{Measurement Constraints}

Beyond missing structural data, several performance-relevant variables are
unavailable in the current dataset. Split times would allow for more specific
analysis of how pool characteristics influence individual phases of a swim,
such as the start, turn, and underwater breakout. Reaction times at the start
block are similarly not captured in meet results as reported by
\textcite{usaswimming2026}. Attendance figures present a related constraint:
while spectator capacity is available as a structural variable, actual
attendance at individual meets is rarely reported, limiting the precision with
which crowd effects can be estimated. These gaps reinforce the AIC finding in
Section~\ref{sec:results} that unmeasured facility dimensions likely account
for a meaningful share of residual venue-level variation.

\subsection{Sample Size Imbalances}

The distribution of observations across facilities is highly uneven. Pools that
host large invitational meets, conference championships, or NCAA postseason
competition contribute substantially more observations to the dataset than pools
used primarily for dual meets or smaller invitationals. This imbalance affects
the precision with which pool-level effects can be estimated: venues with sparse
representation yield wide confidence intervals around their estimated intercepts
and may exercise disproportionate influence on coefficient estimates for
facility-level covariates.

% FIX 4: Expanded Proposed Solutions section from a thin list into
% substantive prose that addresses each limitation directly.
\subsection{Proposed Solutions}

Several strategies are employed to address these limitations within the current
framework. Where direct measurement of a facility characteristic is unavailable,
proxy variables are substituted based on theoretical and practical grounds. Pool
depth is proxied by meet level, on the premise that pools hosting championship
competition are more likely to meet depth standards associated with reduced
turbulence. Gutter design is proxied by facility age, reflecting the historical
adoption of modern wave-dampening gutter technology over the past two decades.
Attendance is proxied by competitive tier---whether NCAA Championship, conference
championship, or standard invitational---which correlates with expected crowd
size even absent direct attendance records.

To address sample size imbalances, the primary analytical sample is restricted
to pools exceeding a minimum observation threshold, as described in
Section~\ref{sec:data}. Robustness checks replicate the primary model on a
restricted subsample of elite facilities---specifically those that have hosted
conference or NCAA championship competition---where data quality and observation
counts are highest. Consistency of estimates across the full and restricted
samples provides evidence that results are not driven by sparse or low-quality
observations at the periphery of the dataset.

Future data collection efforts should prioritize direct measurement of gutter
configuration, lane rope specifications, and water temperature at the time of
competition. Partnering with timing system providers or conference offices to
access split-time archives would substantially enrich the analysis and allow for
a phase-by-phase decomposition of the pool effect, an extension outlined in
Section~\ref{sec:future}.

% ══════════════════════════════════════════════════════════════════════════════
\section{Implications}

\subsection{For Athletes and Coaches}

The findings of this study have direct practical relevance for competitive
swimmers and their coaching staffs. Performance evaluation in NCAA swimming has
historically relied on raw or conversion-adjusted times, with limited systematic
attention to the venue at which a time was posted. The results presented here
suggest that pool characteristics contribute meaningfully to observed performance
variation, implying that times achieved at different facilities are not
straightforwardly comparable without accounting for venue effects. Coaches
interpreting a swimmer's season progression, evaluating recruiting targets, or
setting performance benchmarks would benefit from contextualizing times relative
to the pool in which they were recorded.

The findings also carry implications for strategic meet selection. To the extent
that certain facility characteristics---indoor environment and altitude chief
among them---are associated with faster performances, athletes seeking to
achieve qualifying standards or personal bests may be better positioned to do
so at venues with those characteristics. While factors such as travel cost,
competitive field, and scheduling will always constrain meet selection decisions,
the results suggest that venue quality is a non-trivial input into that calculus.

\subsection{For the NCAA and Athletic Conferences}

At the institutional level, the results raise questions about the criteria
currently used to select sites for conference and NCAA championship competition.
If pool characteristics systematically influence performance, then the choice of
championship venue is not merely a logistical decision but one with competitive
consequences. Facilities that confer measurable speed advantages will tend to
produce faster championship times, affecting qualifying standards, All-American
cutoffs, and the historical record of the sport more broadly.

A related concern involves facility equity. The variation in pool quality
documented in this study is not randomly distributed across programs: larger,
better-resourced athletic departments are more likely to operate well-equipped
natatoriums. If access to fast pools influences performance outcomes, then
institutional resource disparities may compound competitive inequality in ways
that are not currently accounted for in how the sport is administered.

\subsection{For Sports Analytics}

This study contributes to a growing body of quantitative research on contextual
and environmental effects in athletic performance. Venue effects have been
examined extensively in sports such as baseball, soccer, and basketball, but
competitive swimming has received comparatively little attention from the sports
analytics community, despite being one of the most precisely timed and
extensively documented individual sports in the world. The mixed-effects
framework developed here provides a replicable template for future work in
aquatic sports and may generalize to other individual-performance sports where
athletes compete across heterogeneous venues. More broadly, this study
demonstrates that even in a sport defined by hundredths of a second,
environmental and structural context shapes what ends up on the scoreboard.

% ══════════════════════════════════════════════════════════════════════════════
\section{Future Extensions}
\label{sec:future}

The present study establishes a foundation for a broader research program on
venue effects in competitive swimming, and several extensions would
substantially enrich the analysis.

The most immediate methodological refinement would be the incorporation of
individual swimmer random effects alongside team random effects, allowing the
model to leverage within-swimmer variation across pools more directly and to
separate athlete development trajectories from venue influences. A closely
related extension would isolate home pool advantage as a distinct phenomenon,
estimating whether swimmers perform systematically differently at their own
facility relative to comparable away venues, net of the structural
characteristics already modeled. This extension would directly address the
broader home-field advantage literature and connect the present findings to
existing work on venue familiarity effects in individual sports.

Expanding the sample across additional competitive seasons would increase
statistical power, improve the precision of pool-level estimates, and allow for
examination of whether venue effects are stable over time or shift as facilities
are renovated or replaced. A larger sample would also support the inclusion of
a greater number of facilities, reducing the influence of the observation
threshold restrictions currently imposed on the analytical sample.

On the methodological side, hierarchical Bayesian modeling offers a natural
extension of the mixed-effects framework, providing full posterior distributions
over pool-level effects and offering a more principled treatment of uncertainty
in pools with sparse observations. The incorporation of split times and reaction
times, should such data become available through timing system archives or
direct institutional partnerships, would enable a decomposition of the pool
effect across discrete phases of a race.

Finally, extending the framework beyond collegiate competition to Olympic and
elite international swimming would represent a substantively important
generalization. World Aquatics certified pools are nominally standardized, but
variation in depth, gutter design, altitude, and venue atmosphere persists even
at the highest levels of the sport. Applying the analytical approach developed
here to championship and World Cup circuit data would both test the external
validity of the current findings and address questions of direct relevance to
the international swimming community.

% ══════════════════════════════════════════════════════════════════════════════
\newpage
\printbibliography[title={Works Cited}]

\end{document}
